\clearpage\SourceLeaf{116}
\pdfbookmark[1]{Göttingische gelehrte Anzeigen. 2 March 1816.}{anzeige1816}
\begin{center}
\rule{0.45\linewidth}{0.4pt}\par
\textbf{Göttingische gelehrte Anzeigen.\quad 2 March 1816.}\par
\rule{0.45\linewidth}{0.4pt}
\end{center}

After Professor \textsc{Gauss}'s lecture, reported in these Reviews last year, had already been printed, his continued work on the same subject brought him the good fortune of reaching the same goal by an entirely new path. On 30 January he submitted a short paper on this to the Royal Society, entitled:

\begin{center}
\emph{Third proof of the theorem on the factorization of polynomials into real factors.}
\end{center}

The two earlier proofs of the most important theorem in the theory of equations differed in that the first was very short and simple, but rested in part on geometric considerations, whereas the other was purely analytic but much more complicated, so that it was impossible to give an adequate account of it in the limited space of these pages. The present third proof, based on entirely different principles, is likewise purely analytic, but surpasses even the first so greatly in simplicity and brevity that we can quite readily communicate all its essentials here in a few pages.

Let $X$ denote the algebraic function of the indeterminate $x$,
\[
x^m+Ax^{m-1}+Bx^{m-2}+Cx^{m-3}+\mathrm{etc.}
\]
\SourceLeaf{117}
where the coefficients $A,B,C$, etc., are fixed real quantities. Let $r$ and $\varphi$ also be two new indeterminates, and put
\[
\begin{aligned}
&r^m\cos m\varphi+Ar^{m-1}\cos(m-1)\varphi
\\
&\quad+Br^{m-2}\cos(m-2)\varphi+Cr^{m-3}\cos(m-3)\varphi
+\mathrm{etc.}=t\\
&r^m\sin m\varphi+Ar^{m-1}\sin(m-1)\varphi
\\
&\quad+Br^{m-2}\sin(m-2)\varphi+Cr^{m-3}\sin(m-3)\varphi
+\mathrm{etc.}=u\\
&mr^m\cos m\varphi+(m-1)Ar^{m-1}\cos(m-1)\varphi
\\
&\quad+(m-2)Br^{m-2}\cos(m-2)\varphi+(m-3)Cr^{m-3}\cos(m-3)\varphi
+\mathrm{etc.}=t'\\
&mr^m\sin m\varphi+(m-1)Ar^{m-1}\sin(m-1)\varphi
\\
&\quad+(m-2)Br^{m-2}\sin(m-2)\varphi+(m-3)Cr^{m-3}\sin(m-3)\varphi
+\mathrm{etc.}=u'\\
&mmr^m\cos m\varphi+(m-1)^2Ar^{m-1}\cos(m-1)\varphi
\\
&\quad+(m-2)^2Br^{m-2}\cos(m-2)\varphi\\
&\quad+(m-3)^2Cr^{m-3}\cos(m-3)\varphi
+\mathrm{etc.}=t''\\
&mmr^m\sin m\varphi+(m-1)^2Ar^{m-1}\sin(m-1)\varphi
\\
&\quad+(m-2)^2Br^{m-2}\sin(m-2)\varphi\\
&\quad+(m-3)^2Cr^{m-3}\sin(m-3)\varphi
+\mathrm{etc.}=u''\\
&\frac{(tt+uu)(tt''+uu'')+(tu'-ut')^2-(tt'+uu')^2}
{r(tt+uu)^2}=y
\end{aligned}
\]

The factor $r$ can evidently be removed from the denominator of the expression for $y$, since $t',u',t'',u''$ are divisible by $r$.

Further, let $R$ be any fixed positive quantity, subject only to the restriction that it be greater than the greatest of the following:
\[
mA\surd2,\quad\surd(mB\surd2),\quad
{}^3\!\surd(mC\surd2),\quad{}^4\!\surd(mD\surd2)\ \mathrm{etc.}
\]
disregarding the signs of $A,B,C$, etc., or treating them all as positive.

Finally, let the following notation also be used:
\[
\begin{aligned}
&R^m\cos45^0+AR^{m-1}\cos(45^0+\varphi)
\\
&\quad+BR^{m-2}\cos(45^0+2\varphi)+CR^{m-3}\cos(45^0+3\varphi)
+\mathrm{etc.}=T\\
&R^m\sin45^0+AR^{m-1}\sin(45^0+\varphi)
\\
&\quad+BR^{m-2}\sin(45^0+2\varphi)+CR^{m-3}\sin(45^0+3\varphi)
+\mathrm{etc.}=U\\
&mR^m\cos45^0+(m-1)AR^{m-1}\cos(45^0+\varphi)
\\
&\quad+(m-2)BR^{m-2}\cos(45^0+2\varphi)
\\
&\quad+(m-3)CR^{m-3}\cos(45^0+3\varphi)+\mathrm{etc.}=T'\\
&mR^m\sin45^0+(m-1)AR^{m-1}\sin(45^0+\varphi)
\\
&\quad+(m-2)BR^{m-2}\sin(45^0+2\varphi)
\\
&\quad+(m-3)CR^{m-3}\sin(45^0+3\varphi)+\mathrm{etc.}=U'
\end{aligned}
\]

After all these preliminaries, which we have collected together for ease of reference, the following conclusions are readily obtained:

I. The quantity $T$ is necessarily positive, whatever value is assigned to $\varphi$. This is easily seen by putting $T$ into the following form:
\[
\begin{aligned}
&\frac{R^{m-1}}{m\surd2}(R+mA\surd2.\cos(45^{\mathrm{a}}+\varphi))\\
&\quad+\frac{R^{m-2}}{m\surd2}(RR+mB\surd2.\cos(45^0+2\varphi))\\
&\quad+\frac{R^{m-3}}{m\surd2}(R^3+mC\surd2.\cos(45^0+3\varphi))\\
&\quad+\frac{R^{m-4}}{m\surd2}(R^4+mD\surd2.\cos(45^0+4\varphi))\\
&\quad+\mathrm{etc.}
\end{aligned}
\]
\SourceLeaf{118}
since each part, taken separately, is necessarily positive. In a similar way it is easy to prove that for every value of $\varphi$, $U,T',U'$ too are necessarily positive.

II. For $r=R$ the functions $t,u,t',u'$ take the following values, respectively:
\[
\begin{gathered}
T\cos(45^0+m\varphi)+U\sin(45^0+m\varphi)\\
T\sin(45^0+m\varphi)-U\cos(45^0+m\varphi)\\
T'\cos(45^0+m\varphi)+U'\sin(45^0+m\varphi)\\
T'\sin(45^0+m\varphi)-U'\cos(45^0+m\varphi)
\end{gathered}
\]
It follows that for $r=R$, $tt+uu=TT+UU$, $tt'+uu'=TT'+UU'$, so both are necessarily positive.

III. Thus, for no value of $r$ greater than each of the quantities $mA\surd2$, $\surd(mB\surd2)$, ${}^3\!\surd(mC\surd2)$, etc., can $t=0$ and $u=0$ hold simultaneously.

The main theorem is now as follows:

\emph{Within the limits $r=0$, $r=R$, and $\varphi=0$, $\varphi=360^0$ (inclusive), there certainly exist values of $r$ and $\varphi$ that give $t=0$ and $u=0$ simultaneously.}

The proof is given as follows. If the theorem is assumed to be false, it follows that $tt+uu$ must always be positive for all values of $r$ and $\varphi$ between the stated limits, and hence $y$ must always be \emph{finite}. Let us now consider the value of the double integral
\[
\iint y\,\mathrm{d}r\,\mathrm{d}\varphi
\]
taken from $r=0$ to $r=R$ and from $\varphi=0$ to $\varphi=360^0$, which will therefore have a fully determined finite magnitude. This value, to be denoted by $\Omega$, can be found in two ways: by integrating first with respect to $\varphi$ and then with respect to $r$, or in the reverse order. Both ways must necessarily lead to the same result.

Now, however, when $r$ is regarded as constant, the indefinite integral is
\[
\int y\,\mathrm{d}\varphi=\frac{tu'-ut'}{r(tt+uu)}
\]
as can easily be checked by differentiation. No constant quantity
\SourceLeaf{119}
is to be added in the integration, insofar as it is to begin at $\varphi=0$, since for this value evidently
\[
\frac{tu'-ut'}{tt+uu}=0
\]
holds. Since the same holds for $\varphi=360^0$, the value of $\int y\,d\varphi$, taken from $\varphi=0$ to $\varphi=360^0$, is $=0$ for every value of $r$. Thus we conclude that $\Omega=0$.

On the other hand, when $\varphi$ is regarded as constant, the indefinite integral is
\[
\int y\,dr=\frac{tt'+uu'}{tt+uu}
\]
as is likewise easily confirmed by differentiation with respect to $r$. Here too no constant quantity is to be added, insofar as the integration is to begin at $r=0$. The integral $\int y\,dr$, taken from $r=0$ to $r=R$, is therefore, by II,
\[
=\frac{TT'+UU'}{TT+UU}
\]
and hence certainly positive for every value of $\varphi$. It follows that $\Omega$ too, that is, the value of the integral
\[
\int\frac{TT'+UU'}{TT+UU}\,d\varphi
\]
taken from $\varphi=0$ to $\varphi=360^0$, must necessarily be positive, which contradicts the first result. The assumption was therefore inadmissible, and the truth of the theorem itself is thereby proved.

It is well enough known, and need not be developed further here, that for all values of $r$ and $\varphi$ giving $t=0$ and $u=0$ simultaneously,
\[
r(\cos\varphi+\sin\varphi\surd(-1))\quad\hbox{and}\quad
r(\cos\varphi-\sin\varphi\surd(-1))
\]
are roots of the equation $X=0$, or that the function $X$ then contains either the factor of the second degree
\[
xx-2r\cos\varphi.x+rr
\]
(when neither $r$ nor $\sin\varphi$ vanishes), or the factor of the first degree
\[
x\mp r
\]
\SourceLeaf{120}
(when either $\sin\varphi=0$, and hence $\cos\varphi=\pm1$, or $r=0$). What has just been proved of $X$ applies again to the quotient after $X$ has been divided by that factor, and so on; from this the complete factorization of the function $X$ into such factors is evident.

This short sequence of deductions fully achieves the actual goal. An explanation could still be desired, though not required, of how the above contradiction---which was unavoidable if the theorem was regarded as false---disappears when the truth of the theorem is taken as the starting point. The author has given some hints on this, which can only be touched on briefly here. If the theorem is true, not all values of $y$ are finite any longer; it is therefore no longer permissible without further consideration to regard $\iint y\,dr\,d\varphi$ as an actual determined quantity. There is thus no reason to be surprised that the mere blind mechanism of the calculus yields contradictory results by different paths. Analysis very often helps itself in this way: to questions put to it without restriction, even though they may be absurd in certain cases, it gives only partly determined answers. Such is the case in determining the value of the integral $\iint y\,dr\,d\varphi$. If, in general, it is to be taken from $r=k$ to $r=l$ and from $\varphi=K$ to $\varphi=L$, and the value of $\frac{u}{t}$ is denoted
\[
\begin{array}{rcl}
\hbox{for }r=k,\quad\varphi=K&\hbox{by}&\theta\\
r=k,\quad\varphi=L&\hbox{by}&\theta'\\
r=l,\quad\varphi=K&\hbox{by}&\theta''\\
r=l,\quad\varphi=L&\hbox{by}&\theta'''
\end{array}
\]
then the analytic operations give the following expression for that integral:
\[
\mathrm{Arc\,tg}\,\theta-\mathrm{Arc\,tg}\,\theta'
-\mathrm{Arc\,tg}\,\theta''+\mathrm{Arc\,tg}\,\theta'''
\]

The integral in fact has a true value only when $y$ remains finite throughout the stated limits. This true value is indeed included in the expression above, but is not yet fully determined by that expression alone, since the function Arc. tg. is multivalued. Further considerations, which are not difficult, must decide \emph{which}
\SourceLeaf{121}
values of this function are to be taken in the particular case. Whenever, on the other hand, $y$ becomes infinite anywhere within the stated limits, the question of the value of $\iint y\,dr\,d\varphi$ is, strictly speaking, absurd. If, ignoring this absurdity, one nevertheless ventures to answer it, there is no reason to be surprised at obtaining one answer by one path and another by other paths; these different answers are, however, always included in the partly determined expression above.

\par\medskip\noindent\rule{0.25\linewidth}{0.4pt}\par\smallskip
\begingroup\small
\emph{Handwritten remarks.}

The \emph{condition} on $R$ can also be stated by requiring it to be greater than $1$ and than the absolute sum of the quantities $A,B,C$, etc., multiplied by $\surd2$. Let this absolute sum be $=\alpha A+\beta B+\gamma C+\delta D+\mathrm{etc.}$, so that each of the coefficients $\alpha,\beta,\gamma$, etc., is either $+1$ or $-1$. The positive value of $T$, etc., is then clear from the following form:
\[
\begin{aligned}&T=R^{m-1}\cos45^0\bigl(R-\surd2.(\alpha A+\beta B+\gamma C+\mathrm{etc.})\bigr)\\
&\quad+R^{m-1}\bigl(\alpha A+A\cos(45^0+\varphi)\bigr)\\
&\quad+R^{m-2}\bigl(\beta BR+B\cos(45^0+2\varphi)\bigr)+\mathrm{etc.}
\end{aligned}
\]

\emph{Theorem.} If $a,b,c\,.\,.\,.\,m,n$ are the roots of the equation $fx=0$, and $a',b',c'\,.\,.\,.\,m'$ the roots of the equation $f'x=0$, where $f'x=\frac{dfx}{dx}$, and the corresponding points in the plane are denoted by the same letters, then, if equal repelling or attracting masses acting inversely as the distance are imagined at $a,b,c\,.\,.\,.\,m,n$, there is equilibrium at $a',b',c'\,.\,.\,.\,m'$.
\endgroup

\par\medskip\begin{center}\rule{0.25\linewidth}{0.4pt}\end{center}

